\subsubsection{Formulación del problema [Pendientes: 1. delimitar el sector de estudio]}\label{sec:problema}
¿Cómo ofrecer al comandante o coordinador de la operación una estimación consultable y una visualización georreferenciada de la posible propagación de un incendio forestal en un sector de los Cerros Orientales de Bogotá, aprovechando resultados calculados previamente y reduciendo la necesidad de ejecutar simulaciones físicas durante el incidente? La pregunta se concentra en transformar resultados científicos en información accesible para el usuario descrito en la sección~\ref{sec:problematica}, dentro de las condiciones representadas por una biblioteca finita. Por ello, la solución deberá permitir reconocer tanto los escenarios compatibles con la situación observada como sus límites de representación, sin confundir una estimación condicionada por datos y supuestos con una predicción cierta del comportamiento real. El ámbito territorial definitivo permanece pendiente y su selección determinará la extensión de los datos, los casos de simulación y las comprobaciones que podrán sostenerse en el proyecto.
